\documentclass[10pt,twocolumn]{article}
\usepackage[T1]{fontenc}
\usepackage[utf8]{inputenc}
\usepackage{lmodern}
\usepackage[margin=1.8cm,top=2.4cm,bottom=2cm,columnsep=0.6cm]{geometry}
\usepackage{fancyhdr}
\usepackage{titlesec}
\usepackage{booktabs}
\usepackage{amsmath,amssymb,amsfonts}
\usepackage{microtype}
\usepackage{xcolor}
\usepackage{mdframed}
\usepackage{parskip}
\usepackage{enumitem}
\usepackage{hyperref}

\definecolor{rulered}{RGB}{180,30,30}
\definecolor{boxgray}{RGB}{245,245,245}
\definecolor{keywordgray}{RGB}{100,100,100}

\pagestyle{fancy}
\fancyhf{}
\fancyhead[L]{\textsc{\small Claude Mag}}
\fancyhead[C]{\small\textit{ICML 2026 Special Edition: Uncertainty \& Conformal Prediction}}
\fancyhead[R]{\small 2026-07-08}
\fancyfoot[C]{\small\thepage}
\renewcommand{\headrulewidth}{0.8pt}

\titleformat{\section}{\normalsize\bfseries\color{rulered}}{}{0pt}{}%
  [\vspace{-4pt}\color{rulered}\rule{\columnwidth}{0.4pt}]
\titlespacing{\section}{0pt}{8pt}{4pt}

\hypersetup{colorlinks=true,urlcolor=rulered,linkcolor=black}

\begin{document}
\setlength{\parindent}{0pt}
\setlength{\parskip}{4pt}

{\small\color{keywordgray}\textbf{Keywords:}
  conformal risk control \textbullet{}
  robust optimization \textbullet{}
  Pareto frontier \textbullet{}
  predict-then-optimize}\par\vspace{4pt}

{\LARGE\bfseries\raggedright
  Calibrating Decision Robustness via Inverse Conformal Risk
  Control}\par\vspace{4pt}

{\large\itshape\raggedright
  Tracing the Miscoverage-Regret Pareto Frontier Instead of Fixing a
  Coverage Target in Advance}\par\vspace{6pt}

{\small\textbf{By} Wenbin Zhou \& Shixiang Zhu
  (Carnegie Mellon University)}\par\vspace{2pt}

{\small\color{keywordgray}arXiv:2510.07750 \textbullet{}
  ICML 2026 Poster}
\par\vspace{2pt}\noindent{\color{rulered}\rule{\columnwidth}{0.6pt}}\vspace{6pt}

Robust optimization protects decisions by hedging against a conformal
uncertainty set---but the robustness level behind that set is usually
fixed ad hoc, either exposing decisions or wasting performance on
needless conservatism. \textbf{Inverse Conformal Risk Control (ICRC)}
flips the usual recipe: instead of fixing coverage and later measuring
regret, it delivers simultaneous, distribution-free, finite-sample
guarantees on both miscoverage \emph{and} regret across an entire
family of policies, letting decision-makers pick their robustness level
directly off a certified Pareto frontier.

\begin{mdframed}[backgroundcolor=boxgray,linecolor=rulered,linewidth=1pt,
    innerleftmargin=6pt,innerrightmargin=6pt,
    innertopmargin=6pt,innerbottommargin=6pt]
\textbf{\small AT A GLANCE}\par\vspace{4pt}
\begin{description}[leftmargin=0pt,labelsep=4pt,itemsep=2pt,parsep=0pt,
    font=\small\bfseries]
  \item[Problem:] The coverage level behind a robust conformal decision
    policy is normally fixed a priori, with no statistical way to judge
    whether it was the right choice.
  \item[Why it matters:] As conformal sets increasingly feed robust
    optimization pipelines in finance, operations, and safety-critical
    planning, an unexamined robustness knob leaves performance on the
    table invisibly.
  \item[Method:] Simultaneous, distribution-free finite-sample bounds
    on both miscoverage and regret across a family of robustness
    levels, valid uniformly.
  \item[Result:] A certified miscoverage-regret Pareto frontier that
    decision-makers can select from \emph{after} seeing it, with the
    guarantee still holding.
  \item[Evidence:] Robust predict-then-optimize benchmarks comparing
    ICRC's certified frontier against fixed-$\alpha$ conformal robust
    optimization.
\end{description}
\end{mdframed}

\section{The Big Picture}

In robust predict-then-optimize, a conformal uncertainty set
$\mathcal{U}_\alpha(x)$ is built at coverage level $\alpha$, and a
decision minimizes the worst-case cost over that set. The choice of
$\alpha$ determines both the protection the decision receives and the
optimization performance (regret) sacrificed for that protection---but
conventional practice fixes $\alpha$ once, in advance, with no
principled way to know whether that choice was too cautious or not
cautious enough for the task's real cost structure. This paper makes
the entire tradeoff between protection and performance a statistically
certified, inspectable object rather than a single unexamined
hyperparameter.

\section{Why Existing Approaches Fall Short}

Fixing a single coverage target $\alpha$ and only afterward measuring
the regret it induces means the decision-maker learns whether their
choice was good only after committing to it. Re-running the pipeline
for a different $\alpha$ offers no statistical guarantee that the new
choice is valid, since the original calibration was only certified for
the original $\alpha$. This traps practitioners into a single ad hoc
choice with no honest way to compare alternatives.

\section{Core Mechanism}

Rather than treating $\alpha$ as a knob set once, ICRC treats the
entire curve relating miscoverage to regret as the object of interest.
It constructs a family of policies indexed by robustness level and
provides distribution-free, finite-sample bounds on both the true
miscoverage rate and the true regret at every level in that family
\emph{simultaneously}, so a decision-maker can inspect the empirical
Pareto frontier after the fact and select whichever point matches their
actual cost-risk preference---with the guarantee still holding for that
a posteriori choice.

\section{Technical Formulation}

For a family of robust policies $\{\pi_\alpha\}$ built from uncertainty
sets $\mathcal{U}_\alpha(x)$, define
\[
\text{miscoverage}(\alpha) = P\big(y \notin \mathcal{U}_\alpha(x)\big),
\]
\[
\text{regret}(\alpha) = \mathbb{E}\big[\text{cost}(\pi_\alpha(x), y)\big]
  - \min_{a} \mathbb{E}[\text{cost}(a, y)].
\]
These move in opposite directions as $\alpha$ varies, forming a Pareto
frontier. ICRC constructs simultaneous confidence bands for both
functions, valid uniformly over a grid of $\alpha$ values with
probability at least $1-\delta$, using conformal calibration on a
held-out set combined with uniform concentration arguments (in the
style of DKW-type bounds) to control the multiple-testing-like
inflation from evaluating many $\alpha$ values on the same data.

\section{Learning or Inference Procedure}

\begin{enumerate}[itemsep=2pt,parsep=0pt]
  \item Build a grid of candidate robustness levels $\alpha$ and their
    corresponding uncertainty sets $\mathcal{U}_\alpha(x)$.
  \item On held-out calibration data, estimate miscoverage and regret
    for every $\alpha$ in the grid.
  \item Apply uniform concentration bounds to certify simultaneous
    validity of both estimates across the whole grid.
  \item Present the certified empirical Pareto frontier; the
    decision-maker selects $\hat\alpha$ off the frontier according to
    their cost-risk preference.
\end{enumerate}

\section{What the Guarantee Says}

With probability at least $1-\delta$, the true miscoverage and true
regret curves are simultaneously bracketed by the certified bands at
every $\alpha$ on the grid---so whichever $\hat\alpha$ the
decision-maker picks \emph{after} inspecting the frontier still enjoys
valid finite-sample guarantees on both quantities, unlike a guarantee
calibrated for one $\alpha$ chosen in advance.

\section{Experimental Findings}

On robust predict-then-optimize benchmarks, ICRC's certified Pareto
frontier lets practitioners identify robustness levels that achieve
materially lower regret than conventional fixed-$\alpha$ conformal
robust optimization at matched miscoverage, and vice versa---revealing
that ad hoc $\alpha$ choices in prior pipelines were leaving performance
on the table without practitioners ever being able to see it.

\section{Ablations and Interpretation}

\textbf{Grid resolution:} Finer grids over $\alpha$ tighten the
achievable frontier at a modest cost in the width of the simultaneous
confidence bands, reflecting the usual multiple-testing tradeoff.

\textbf{Comparison to fixed-$\alpha$ baselines:} Any single fixed-$\alpha$
policy corresponds to one point on ICRC's frontier; the benefit is
purely in being able to see and choose among the alternatives with
statistical validity.

\textbf{Practical takeaway:} ICRC does not require the analyst to
predict the right robustness level in advance---it makes the tradeoff
observable and lets the choice come last.

\section*{Reference}

Wenbin Zhou, Shixiang Zhu.
\textbf{Calibrating Decision Robustness via Inverse Conformal Risk
Control}.
arXiv:2510.07750, October 2025.
\url{https://arxiv.org/abs/2510.07750}

\end{document}
